% The PDF template of the report (STYLE 10.9: no overfull line over 10 pt). Three parts: inline code may
% break after _ / : and @; running heads show the chapter title; long hex strings may break inside.
%
% Inline code may break after _ / : and @.
% Pandoc writes inline code as \texttt{...}. This wraps \texttt so that every \_, /, : and @ in its
% argument is followed by an empty discretionary: a line break with no hyphen, at the penalty
% \exhyphenpenalty, taken only where the line needs it. A C++ scope :: stays together and breaks
% after its second colon. Prose, headings and code blocks (Shaded, verbatim) do not change: blocks do
% not use \texttt. PDF bookmarks do not change: hyperref's \pdfstringdef replaces \texttt by its
% argument. The patterns are strings (catcode 12), as the characters are in the document body,
% whatever \ExplSyntaxOn or \makeatletter make of : and @ here.
\ExplSyntaxOn
\NewCommandCopy \__snapy_texttt:n \texttt
\cs_new:Npn \__snapy_tt_break: { \discretionary { } { } { } }
\cs_generate_variant:Nn \tl_replace_all:Nnn { Nee }
\tl_new:N \l__snapy_tt_tl
\cs_new_protected:Npn \__snapy_tt_after:n #1
  {
    \tl_replace_all:Nee \l__snapy_tt_tl
      { \tl_to_str:n {#1} } { \tl_to_str:n {#1} \exp_not:N \__snapy_tt_break: }
  }
\RenewDocumentCommand \texttt { m }
  {
    \tl_set:Nn \l__snapy_tt_tl {#1}
    \tl_replace_all:Nnn \l__snapy_tt_tl { \_ } { \_ \__snapy_tt_break: }
    \__snapy_tt_after:n { / }
    \__snapy_tt_after:n { @ }
    \__snapy_tt_after:n { : }
    \tl_replace_all:Nee \l__snapy_tt_tl
      { \exp_not:N \__snapy_tt_break: \tl_to_str:n { : } } { \tl_to_str:n { : } }
    \__snapy_texttt:n { \tl_use:N \l__snapy_tt_tl }
  }
\ExplSyntaxOff

% Running heads show the chapter title only, on left and right pages (editor's ruling): a section title
% with a long code name does not fit a one-line head, and authors do not shorten their titles.
\usepackage{scrlayer-scrpage}
\automark[chapter]{chapter}
\pagestyle{scrheadings}

% Long hex strings in running text (a 40-character sha, path@sha) may break anywhere inside, with no
% hyphen (STYLE 10.9). A Lua filter, run on each paragraph before line breaking, puts an empty
% discretionary (penalty \exhyphenpenalty) between the characters of every run of at least 16
% lower-case hex characters that holds a digit and a letter, outside math. Ordinary words never
% qualify, so their hyphenation is unchanged.
\directlua{
  local GLYPH, KERN, DISC, MATH = node.id("glyph"), node.id("kern"), node.id("disc"), node.id("math")
  local function hex(c) return (c >= 48 and c <= 57) or (c >= 97 and c <= 102) end
  local function hexbreaks(head)
    local n, inmath = head, false
    while n do
      if n.id == MATH then
        inmath = (n.subtype == 0)
        n = n.next
      elseif inmath or n.id ~= GLYPH or not hex(n.char) then
        n = n.next
      else
        local run, k, digit, letter, m = {}, 0, false, false, n
        while m and ((m.id == GLYPH and hex(m.char)) or m.id == KERN) do
          if m.id == GLYPH then
            k = k + 1
            run[k] = m
            if m.char <= 57 then digit = true else letter = true end
          end
          m = m.next
        end
        if k >= 16 and digit and letter then
          for i = 1, k - 1 do
            local d = node.new(DISC)
            d.penalty = tex.exhyphenpenalty
            head = node.insert_after(head, run[i], d)
          end
        end
        n = m
      end
    end
    return head
  end
  luatexbase.add_to_callback("pre_linebreak_filter", hexbreaks, "snapy-report hex breaks")
}
